\chapter{Tests et résultats}
\label{chap:tests}

\section{Méthodologie}

Le projet ne dispose pas pour l'instant d'un \emph{framework} de
tests unitaires automatisés~: la nature visuelle de la
simulation et le couplage fort avec GTK rendent peu rentable la
mise en place d'un harnais comme \code{cmocka} ou
\code{check}. À la place, la validation s'est appuyée sur~:

\begin{itemize}
  \item \textbf{Tests fonctionnels manuels} reproductibles,
    listés dans la section suivante. Chaque scénario est défini
    par~: une condition initiale (état du fichier de
    configuration), une action utilisateur, un résultat attendu,
    et un résultat observé.
  \item \textbf{Tests de chargement de configuration} avec des
    cas volontairement erronés~: typo dans une clé, identifiant
    de proie inconnu, fichier introuvable.
  \item \textbf{Compilation avec drapeaux stricts}~:
    \code{-Wall -Wextra -Wpedantic} pour détecter les
    constructions douteuses.
\end{itemize}

Cette approche reste qualitative~: aucune mesure quantitative de
performance (FPS moyen, latence d'un tick) n'a été instrumentée,
ce qui constitue une limite identifiée et reportée dans les
perspectives (chapitre~\ref{chap:conclusion}).

\section{Scénarios fonctionnels}

\subsection{S1 -- Démarrage avec configuration par défaut}

\begin{description}[leftmargin=3em]
  \item[Préconditions.] Le binaire \file{build/bin/aquarium} est
    construit. Le fichier \file{aquarium.conf} contient les
    huit espèces de la configuration de référence avec leurs
    valeurs documentées.
  \item[Action.] Lancer \code{./build/bin/aquarium}.
  \item[Attendu.] La fenêtre s'ouvre en 1280$\times$720, la
    barre latérale présente huit cases à cocher (une par
    espèce), la combobox \emph{Add precise} liste les huit
    noms, le pool initial spawn environ 95 entités au total
    (45 goldfish, 18 clownfish, 14 angelfish, 10 par prédateur
    intermédiaire, 4 tuna, 2 shark).
  \item[Observé.] Conforme. Les bancs se forment en quelques
    secondes, les prédateurs commencent à chasser après le
    premier croisement avec une proie.
\end{description}

\subsection{S2 -- Modification d'un paramètre simple}

\begin{description}[leftmargin=3em]
  \item[Préconditions.] Application en cours d'exécution.
  \item[Action.] Éditer \file{aquarium.conf}, remplacer
    \code{max\_speed = 85.0} par \code{max\_speed = 200.0} dans
    le bloc \code{[[species]]} \code{goldfish}, sauvegarder, et
    cliquer sur \emph{Reload config}.
  \item[Attendu.] Les goldfish se déplacent visiblement plus
    vite (vitesse maximale environ 2,4$\times$ supérieure). Les
    autres espèces conservent leur comportement nominal.
  \item[Observé.] Conforme. Les goldfish atteignent visuellement
    une vitesse plus élevée. Le banc paraît plus dispersé car la
    cohésion à pondération inchangée a moins de prise sur des
    individus rapides.
\end{description}

\subsection{S3 -- Suppression d'une espèce}

\begin{description}[leftmargin=3em]
  \item[Préconditions.] Application en cours d'exécution.
  \item[Action.] Commenter intégralement le bloc
    \code{[[species]]} \code{angelfish} dans la
    configuration en préfixant chaque ligne par \code{\#},
    puis cliquer sur \emph{Reload config}.
  \item[Attendu.] L'espèce angelfish disparaît~: aucun spawn,
    plus de case à cocher, plus d'entrée dans la combobox.
    Les autres espèces continuent normalement. Aucun crash.
  \item[Observé.] Conforme. Les références \code{prey}
    pointaient déjà vers angelfish dans la configuration du
    requin~: l'absence d'angelfish est correctement gérée
    car le bit correspondant n'est simplement pas mis dans
    \code{prey\_mask}.
  \item[Note.] Si une espèce dont l'identifiant figure encore
    dans une liste \code{prey} d'un autre prédateur est
    supprimée, le chargement échoue avec un message clair (cf.\
    scénario~S5).
\end{description}

\subsection{S4 -- Ajout d'une nouvelle espèce}

\begin{description}[leftmargin=3em]
  \item[Préconditions.] Un fichier PNG \file{piranha.png}
    orienté vers la gauche est placé dans \file{assets/}.
  \item[Action.] Ajouter à la fin de \file{aquarium.conf} un
    bloc \code{[[species]]} avec \code{id = "piranha"} et un
    sous-ensemble cohérent de paramètres (vitesse moyenne,
    rayon de chasse, \code{prey = ["goldfish"]}). Cliquer sur
    \emph{Reload config}.
  \item[Attendu.] Une neuvième case à cocher apparaît, une
    neuvième entrée dans la combobox, et des piranhas spawnent
    dans la simulation.
  \item[Observé.] Conforme. Aucun changement de code C n'a été
    nécessaire.
\end{description}

\subsection{S5 -- Configuration invalide : référence inconnue}

\begin{description}[leftmargin=3em]
  \item[Préconditions.] Application en cours d'exécution.
  \item[Action.] Modifier \code{prey = ["goldfish"]} en
    \code{prey = ["sardine"]} dans le bloc \code{trout}, puis
    cliquer sur \emph{Reload config}.
  \item[Attendu.] La boîte de dialogue d'erreur s'affiche avec
    un message du type \og species 'trout': unknown prey
    'sardine'~\fg. La simulation courante n'est pas affectée
    (rollback implicite~: l'ancienne configuration reste en
    place).
  \item[Observé.] Conforme.
\end{description}

\subsection{S6 -- Configuration invalide : syntaxe TOML}

\begin{description}[leftmargin=3em]
  \item[Préconditions.] Application en cours d'exécution.
  \item[Action.] Introduire volontairement une erreur de
    syntaxe (par exemple~: oubli des guillemets autour d'une
    chaîne, ou crochet ouvert non refermé) puis cliquer sur
    \emph{Reload config}.
  \item[Attendu.] La boîte d'erreur affiche le chemin du
    fichier suivi du message d'erreur retourné par
    \code{tomlc99}, qui inclut typiquement le numéro de
    ligne.
  \item[Observé.] Conforme. Le message est précis car
    \code{tomlc99} fournit lui-même un \code{errbuf} explicite.
\end{description}

\subsection{S7 -- Absence du fichier}

\begin{description}[leftmargin=3em]
  \item[Préconditions.] Renommer temporairement
    \file{aquarium.conf} en \file{aquarium.conf.bak} et
    s'assurer qu'aucune copie de \emph{fallback} n'existe en
    \file{\textasciitilde/.config/gtkaqua/} ni en
    \file{/etc/gtkaqua/}.
  \item[Action.] Lancer l'application.
  \item[Attendu.] Échec propre avec le message \og no
    aquarium.conf found (checked ./aquarium.conf,
    \textasciitilde/.config/gtkaqua/aquarium.conf,
    /etc/gtkaqua/aquarium.conf)~\fg. Pas de \emph{segfault}.
  \item[Observé.] Conforme.
\end{description}

\subsection{S8 -- Fallback XDG}

\begin{description}[leftmargin=3em]
  \item[Préconditions.] \file{aquarium.conf} retiré du
    répertoire courant. Une copie placée dans
    \file{\textasciitilde/.config/gtkaqua/aquarium.conf} via
    \code{make install-conf}.
  \item[Action.] Lancer l'application depuis un répertoire
    arbitraire qui ne contient pas \file{aquarium.conf}.
  \item[Attendu.] L'application démarre normalement, en
    utilisant la copie XDG.
  \item[Observé.] Conforme.
\end{description}

\subsection{S9 -- Pause et reprise}

\begin{description}[leftmargin=3em]
  \item[Préconditions.] Simulation en cours.
  \item[Action.] Cliquer sur \emph{Pause}, attendre quelques
    secondes, observer le comportement, puis cliquer sur
    \emph{Resume}.
  \item[Attendu.] Pendant la pause, les sprites restent
    immobiles, les widgets sont à leur dernière position. À la
    reprise, le mouvement reprend sans saut.
  \item[Observé.] Conforme. \code{world\_tick} n'est plus
    appelé pendant la pause, mais le timer GTK continue, donc
    \code{elapsed} ne progresse pas.
\end{description}

\subsection{S10 -- Plein écran}

\begin{description}[leftmargin=3em]
  \item[Préconditions.] Simulation en cours.
  \item[Action.] Appuyer sur \kw{F11}.
  \item[Attendu.] La fenêtre passe en plein écran. Les sprites
    sont libres d'occuper toute la zone, qui est reconfigurée
    automatiquement (\code{world.width}, \code{world.height}
    sont mises à jour à partir de
    \code{gtk\_widget\_get\_width/height} dans le \code{tick\_cb}).
  \item[Observé.] Conforme. Appuyer sur \kw{Échap} restaure le
    mode fenêtré.
\end{description}

\section{Synthèse des résultats}

Le tableau~\ref{tab:test-summary} récapitule les dix scénarios.

\begin{table}[h!]
\centering
\footnotesize
\begin{tabular}{@{}lll@{}}
\toprule
\textbf{Scénario} & \textbf{Fonctionnalité testée} & \textbf{Résultat} \\
\midrule
S1  & Démarrage avec config par défaut       & OK \\
S2  & Modification d'un paramètre simple     & OK \\
S3  & Suppression d'une espèce               & OK \\
S4  & Ajout d'une nouvelle espèce            & OK \\
S5  & Référence \code{prey} inconnue         & OK \\
S6  & Erreur de syntaxe TOML                 & OK \\
S7  & Absence totale du fichier              & OK \\
S8  & Fallback XDG                           & OK \\
S9  & Pause/Resume                           & OK \\
S10 & Plein écran                            & OK \\
\bottomrule
\end{tabular}
\caption{Synthèse des scénarios de test.}
\label{tab:test-summary}
\end{table}

\section{Limites identifiées}

\begin{itemize}
  \item \textbf{Pas de mesure quantitative de performance.} Le
    coût d'un tick n'est pas instrumenté. Pour une évaluation
    rigoureuse, il faudrait mesurer la durée moyenne et le 99e
    centile de \func{world\_tick} en fonction du nombre
    d'entités, et tracer le résultat.

  \item \textbf{Complexité $O(N^2)$.} Toutes les forces
    parcourent la liste complète des entités à chaque tick. À
    population modeste (~100 entités) cela ne pose pas de
    problème, mais le coût croît rapidement~: 200 entités
    impliquent quatre fois plus de comparaisons.

  \item \textbf{Pas de tests automatisés.} Les régressions sont
    actuellement détectées \emph{a posteriori} par observation
    visuelle. Un harnais de tests sur les fonctions pures
    (parsing TOML, résolution de \code{prey\_mask},
    \func{vec2\_*}) serait peu coûteux à mettre en place et
    procurerait une garantie utile.

  \item \textbf{Couverture du parseur de configuration.} Toutes
    les clés sont optionnelles côté parseur~: si une clé
    obligatoire est manquante (par exemple \code{sprite}), le
    chargement détecte le cas, mais d'autres clés silencieusement
    par défaut à zéro pourraient produire des comportements
    dégénérés (vitesse nulle, rayons nuls). Une validation
    \emph{post-load} renforcée serait souhaitable.
\end{itemize}
